\Article{1}
Die Kraft, die einer in ihrem Schwerpunkt aufgehängten Magnetnadel an jedem Ort der Erde eine bestimmte Richtung gibt, nennt man die erdmagnetische Kraft, insofern der Sitz ihrer Ursache nur im Erdkörper selbst zu suchen ist. Dabei setzt man voraus, dass jede fremde äußere Ursache, die auf die Nadel wirken könnte, beseitigt ist, etwa die Nähe eines anderen künstlichen Magneten oder die Nähe des Leiters eines galvanischen Stroms. Allerdings ist fraglich, ob die regelmäßigen und unregelmäßigen stündlichen Änderungen dieser Kraft ihre unmittelbaren Ursachen nicht außerhalb des Erdkörpers haben könnten. Es ist zu hoffen, dass die Aufmerksamkeit, die die Naturforscher diesen Erscheinungen jetzt allgemein widmen, uns darüber künftig bedeutende Aufschlüsse geben wird. Man darf aber nicht vergessen, dass diese Änderungen vergleichsweise sehr klein sind. Es muss also eine viel stärkere, beständig wirkende Hauptkraft geben, deren Sitz wir in der Erde selbst annehmen. Daraus folgt unmittelbar, dass man die tatsächliche Grundlage für die Untersuchung dieser Hauptkraft eigentlich von den erwähnten anomalen Änderungen befreien sollte. Das ist nur durch \SourceLeaf{130}Mittelwerte aus zahlreichen fortgesetzten Beobachtungen möglich. Solange solche von den Änderungen bereinigten Ergebnisse für eine große Zahl von Punkten auf der gesamten Erdoberfläche fehlen, lässt sich höchstens eine Annäherung erreichen, bei der Abweichungen von der Größenordnung dieser Anomalien zurückbleiben können.

\Article{2}
Unseren Untersuchungen liegt die Annahme zugrunde, dass die erdmagnetische Kraft die Gesamtwirkung der magnetisierten Teile des Erdkörpers ist. Wir stellen uns die Magnetisierung als Trennung der magnetischen Flüssigkeit vor. Hat man diese Vorstellung einmal angenommen, dann gehört die Wirkungsweise dieser Flüssigkeiten zu den erwiesenen physikalischen Tatsachen. Gleichnamiges stößt sich ab, Ungleichnamiges zieht sich an, jeweils im umgekehrten Verhältnis zum Quadrat der Entfernung. Würden wir diese Vorstellung durch die von \textsc{Ampère} ersetzen, nach der es keine magnetischen Flüssigkeiten gibt und der Magnetismus allein in beständigen galvanischen Strömen in den kleinsten Teilen der Körper besteht, so änderte sich an den Ergebnissen gar nichts. Dasselbe würde gelten, wenn wir den Erdmagnetismus einer gemischten Ursache zuschrieben, sodass er teils aus der Trennung der magnetischen Flüssigkeiten in der Erde, teils aus galvanischen Strömen in ihr hervorginge. Bekanntlich kann man nämlich an die Stelle eines jeden galvanischen Stroms eine bestimmte Verteilung magnetischer Flüssigkeiten auf einer von der Stromlinie begrenzten Fläche setzen, die in jedem Punkt des äußeren Raums genau dieselbe magnetische Wirkung ausübt wie der galvanische Strom.

\Article{3}
Zur Messung der magnetischen Flüssigkeiten nehmen wir wie in der Schrift \textit{Intensitas vis magneticae} usw. als positive Einheit diejenige Menge nördlicher Flüssigkeit, die auf eine ebenso große Menge derselben Flüssigkeit in der als Einheit angenommenen Entfernung eine bewegende Kraft ausübt, die gleich der als Einheit angenommenen Kraft ist. Wenn wir einfach von der magnetischen Kraft sprechen, die an einem beliebigen Punkt des Raums durch magnetische Flüssigkeit an einem anderen Ort hervorgerufen wird, verstehen wir darunter stets \SourceLeaf{131}die bewegende Kraft, die dort auf eine Einheit der positiven magnetischen Flüssigkeit ausgeübt wird. In diesem Sinne übt also die in einem Punkt konzentriert gedachte magnetische Flüssigkeit $\mu$ in der Entfernung $\rho$ die magnetische Kraft $\frac{\mu}{\rho\rho}$ aus, und zwar abstoßend oder anziehend in Richtung der geraden Linie $\rho$, je nachdem, ob $\mu$ positiv oder negativ ist. Bezeichnen wir mit $a,b,c$ die Koordinaten von $\mu$ in Bezug auf drei einander rechtwinklig schneidende Achsen und mit $x,y,z$ die Koordinaten des Punkts, an dem die Kraft ausgeübt wird, so gilt $\rho=\BareRoot{((x-a)^2+(y-b)^2+(z-c)^2)}$. Zerlegen wir die Kraft parallel zu den Koordinatenachsen, so sind ihre Komponenten
\[
\frac{\mu(x-a)}{\rho^3},\qquad \frac{\mu(y-b)}{\rho^3},\qquad \frac{\mu(z-c)}{\rho^3}
\]
die, wie man leicht sieht, den partiellen Ableitungen von $-\frac{\mu}{\rho}$ nach $x,y$ und $z$ gleich sind.

Wirken außer $\mu$ noch weitere Teile magnetischer Flüssigkeit, $\mu',\mu'',\mu'''$ usw., die in Punkten konzentriert sind, deren Entfernungen vom Wirkungspunkt jeweils $\rho',\rho'',\rho'''$ usw. betragen, so sind die Komponenten der gesamten daraus entstehenden magnetischen Kraft parallel zu den Koordinatenachsen gleich den partiellen Ableitungen von
\[
-\left(\frac{\mu}{\rho}+\frac{\mu'}{\rho'}+\frac{\mu''}{\rho''}+\frac{\mu'''}{\rho'''}+\text{u. s. w.}\right)
\]
nach $x,y$ und $z$.

